% English translation of checked_pages/esq_119_014.tex; source: cote 119, site 14, PDF page 15, typescript page 12.
% AI-translated by GPT-6 Astra (OpenAI, ChatGPT Pro mode).
\scanpage{119/14}
\mspage{12}

line is defined over the absolute base field $\mathbb{Q}$ ,
as are the allowed ramification points, the algebraic curves
obtained are defined not only over $\mathbb{C}$, but over the algebraic
closure $\overline{\mathbb{Q}}$ of $\mathbb{Q}$ in $\mathbb{C}$. As for the original map, it
is recovered on the algebraic curve as the inverse image of the real
segment $[0,1]$ (where $0$ is regarded as a vertex, and $1$ as
the midpoint of a ``folded edge'' with 1 as its center), which constitutes
the ``universal oriented 2-map'' on the Riemann sphere (*).
The points of the algebraic curve $X$ above 0, 1 and $\infty$
are precisely the vertices and the ``centers'' of the edges and
faces respectively of the map $(X,K)$, and the orders of the vertices
and faces are precisely the multiplicities of the zeros and
poles of the rational function (defined over $\overline{\mathbb{Q}}$) on $X$,
expressing its structural projection to $\mathbb{P}^1_{\mathbb{C}}$.

This discovery, which technically amounts to so little,
made a very strong impression on me, and represents
a decisive turning point in the course of my reflections, in particular
a shift in my center of interest in mathematics, which
suddenly became sharply localized. I do not believe that a
mathematical fact has ever struck me as much as this one, or had
a comparable psychological impact ($^2$). This surely has to do with the
so familiar, nontechnical nature of the objects under consideration,
of which any child's drawing scribbled on a scrap of paper (provided
the drawing is all in one piece) gives a perfectly
explicit example. With such a drawing are associated subtle arithmetic
invariants, which will be completely upset as soon as
one adds another stroke. Since we are dealing here with spherical maps, necessarily

\noindent\rule{4cm}{0.4pt}

(*) There is an analogous description of finite unoriented maps,
possibly with boundary, in terms of \uline{real} algebraic curves,
more precisely of a covering of $\mathbb{P}^1_{\mathbb{R}}$ ramified only at 0, 1, $\infty$\frapparatus{V-FR-014-001},
the surface with boundary associated with such a covering being $X(\mathbb{C})/\tau$, where $\tau$
is complex conjugation. The ``universal'' unoriented map is
here the disk, or upper hemisphere of the Riemann sphere, equipped,
as before, with the embedded 1-complex $K = [0,1]$.

\margin{14}

